\documentclass[conference]{IEEEtran}

\usepackage{graphicx}
\usepackage{amsmath,amssymb}
\usepackage{float}
\usepackage{booktabs}
\usepackage{multirow}

\begin{document}

\title{
    \textbf{Benchmark Report: 525.x264\_r} \\
    \text{Register Spill Analysis on RISC-V via gem5}
}

\author{
    \IEEEauthorblockN{Lütfullah Burak Kaya}
    \IEEEauthorblockA{
        Ozyegin University \\
        burak.kaya.50711@ozu.edu.tr
    }
    \and
    \IEEEauthorblockN{Ismail Akturk}
    \IEEEauthorblockA{
        Ozyegin University \\
        ismail.akturk@ozyegin.edu.tr
    }
}

\newcommand{\LongDate}{%
    \number\day\space
    \ifcase\month
    \or January\or February\or March\or April\or May\or June%
    \or July\or August\or September\or October\or November\or December%
    \fi
    \space \number\year
}

\twocolumn[
\begin{@twocolumnfalse}
\maketitle
\begin{center}{\small \LongDate}\end{center}
\vspace{1em}
\end{@twocolumnfalse}
]

% ================================================

\begin{abstract}
This report presents the register spill characterization of the
\texttt{525.x264\_r} benchmark from SPEC CPU2017, executed on a
RISC-V (RV64GC) target using the gem5 TimingSimpleCPU simulator.
A custom SpillDetector was integrated into gem5 to count store-then-load
patterns on the stack within a user-defined Region of Interest (ROI).
The \texttt{ldecod\_r} binary (JM reference H.264 decoder) was used,
with \texttt{BuckBunny.264} as the test input. Results show a spill
rate of \textbf{3.28\%} within the ROI, corresponding to 768 million
detected spill events across 23.4 billion ROI instructions. H.264
block-based decoding alternates between highly regular CABAC entropy
decoding and IDCT/motion-compensation loops, producing moderate and
relatively consistent register pressure throughout the decode pipeline.
\end{abstract}

% =========================================================
\section{Benchmark Overview}

\texttt{525.x264\_r} is the SPEC CPU2017 rate variant of the x264
H.264/AVC video codec suite. The SPEC benchmark exercises two
separate workloads:
\begin{itemize}
    \item \textbf{x264} (encoder): encodes a YUV video sequence to
          an H.264 bitstream.
    \item \textbf{ldecod\_r} (decoder): decodes an H.264 bitstream to
          a raw YUV file using the JM reference software decoder.
\end{itemize}
This study uses the \textbf{ldecod\_r} binary, which implements the
full H.264 syntax parsing, CABAC/CAVLC entropy decoding, inverse
discrete cosine transform (IDCT), motion compensation, and in-loop
deblocking filter. The decoder processes the bitstream slice by slice,
maintaining the decoded picture buffer (DPB) and reference frame lists.

The RISC-V binary (\texttt{ldecod\_r.riscv}) was compiled with gem5
ROI markers inserted in \texttt{src/ldecod\_src/decoder\_test.c},
wrapping the main per-frame decode loop:

\begin{itemize}
    \item \texttt{m5\_work\_begin(0,0)} --- before the do-while decode loop
    \item \texttt{m5\_work\_end(0,0)}   --- after all frames are decoded
\end{itemize}

This ROI encompasses the complete decoding pipeline for every frame
in the input bitstream.

\subsection{Input Datasets}

SPEC CPU2017 provides three dataset sizes for \texttt{x264\_r},
differing in video content, resolution, and the number of frames
encoded and decoded (Table~\ref{tab:inputs}).

\begin{table}[H]
\centering
\caption{x264\_r (ldecod\_r) Input Dataset Sizes}
\label{tab:inputs}
\begin{tabular}{lll}
\toprule
\textbf{Dataset} & \textbf{Input} & \textbf{Description} \\
\midrule
test    & BuckBunny.264 (2.4\,MiB)  & Short H.264 clip \\
train   & BuckBunny.264 (larger)    & Extended clip \\
refrate & BuckBunny.264 (full)      & Full-length clip \\
\bottomrule
\end{tabular}
\end{table}

The \textbf{test} dataset was selected for this study.
\texttt{BuckBunny.264} is a short segment of the Big Buck Bunny
open-source animated film, encoded in H.264/AVC. Despite the
compact file size (2.4\,MiB), the JM decoder is computationally
intensive: it performs full compliance-level syntax checking, CABAC
context initialisation for each slice, 4$\times$4 and 8$\times$8
IDCT on every macroblock, and motion compensation with sub-pixel
interpolation. The 23.4 billion ROI instructions reflect the
instruction-level cost of a full-compliance software H.264 decoder.

% =========================================================
\section{Simulation Setup}

\subsection{gem5 Configuration}

\begin{itemize}
    \item \textbf{CPU model:}   TimingSimpleCPU
    \item \textbf{ISA:}         RISC-V RV64GC
    \item \textbf{Caches:}      enabled (default L1 I/D + L2)
    \item \textbf{Memory:}      4\,GB simulated physical RAM
    \item \textbf{Spill mode:}  summary only (\texttt{spill\_verbose=False})
\end{itemize}

The full gem5 invocation was:

\begin{verbatim}
./build/RISCV/gem5.opt
  --outdir=benchmarks/x264/m5out_test
  configs/deprecated/example/se.py
  --cmd=benchmarks/x264/ldecod_r.riscv
  --options="-i benchmarks/x264/BuckBunny.264
             -o benchmarks/x264/BuckBunny.yuv"
  --cpu-type=TimingSimpleCPU
  --caches
  --mem-size=4GB
\end{verbatim}

The \texttt{-i} flag specifies the input H.264 bitstream and
\texttt{-o} specifies the output raw YUV file. The JM decoder
reads the input file sequentially, processing each NAL unit in
display order.

% =========================================================
\section{Simulation Results}

The simulation completed successfully. The JM decoder processed
all frames in \texttt{BuckBunny.264} and wrote the decoded output
to \texttt{BuckBunny.yuv}.

\subsection{Pre-ROI Context}

Before the decode loop ROI, the binary executed JM decoder
initialisation: parsing the decoder configuration, allocating
the decoded picture buffer, and initialising CABAC context tables.
Table~\ref{tab:preroi} shows the global counters at ROI entry.

\begin{table}[H]
\centering
\caption{Global Counters at ROI Entry}
\label{tab:preroi}
\begin{tabular}{lr}
\toprule
\textbf{Metric} & \textbf{Value} \\
\midrule
Total instructions (pre-ROI) & 268,194 \\
Total spills detected (pre-ROI) & 8,865 \\
Pre-ROI spill rate & 3.31\% \\
\bottomrule
\end{tabular}
\end{table}

The pre-ROI phase is exceptionally small: only 268,194 instructions,
making it the second smallest startup overhead of all evaluated
benchmarks. The JM decoder initialises relatively little global
state before entering the decode loop, with most data structures
allocated lazily at first-frame decode time inside the ROI.

\subsection{ROI Spill Statistics}

Table~\ref{tab:roi} presents the spill statistics collected within
the ROI.

\begin{table}[H]
\centering
\caption{ROI Spill Statistics --- 525.x264\_r / ldecod\_r (BuckBunny.264)}
\label{tab:roi}
\begin{tabular}{lr}
\toprule
\textbf{Metric} & \textbf{Value} \\
\midrule
ROI instructions         & 23,405,758,981 \\
ROI stores               &  2,401,493,613 \\
ROI loads                &  4,837,558,326 \\
\midrule
\textbf{ROI spills}      & \textbf{767,981,981} \\
\textbf{Spill rate}      & \textbf{3.2812\%} \\
\bottomrule
\end{tabular}
\end{table}

\subsection{Timing}

\begin{table}[H]
\centering
\caption{Simulation Timing}
\label{tab:timing}
\begin{tabular}{lr}
\toprule
\textbf{Metric} & \textbf{Value} \\
\midrule
Simulated time           & 45.63\,s \\
Host wall time           & 16{,}720.87\,s ($\approx$4.6\,hours) \\
Total simulated instructions & 23,406,228,013 \\
\bottomrule
\end{tabular}
\end{table}

% =========================================================
\section{Analysis}

\subsection{Spill Rate Interpretation}

The measured spill rate of \textbf{3.28\%} places \texttt{x264\_r}
eleventh among all evaluated SPEC CPU2017 benchmarks
(Table~\ref{tab:compare}), in the moderate-pressure tier between
\texttt{531.deepsjeng\_r} (4.17\%) and \texttt{544.nab\_r} (2.26\%).

\begin{table}[H]
\centering
\caption{Cross-Benchmark Spill Rate Comparison}
\label{tab:compare}
\begin{tabular}{lrr}
\toprule
\textbf{Benchmark} & \textbf{ROI Spills} & \textbf{Spill Rate} \\
\midrule
507.cactuBSSN\_r & 3,782,099,438 & 7.8756\% \\
511.povray\_r    &   174,359,922 & 7.7809\% \\
500.perlbench\_r &     1,070,055 & 7.2838\% \\
520.omnetpp\_r   &   894,619,681 & 7.1994\% \\
502.gcc\_r       &     1,108,838 & 6.9185\% \\
526.blender\_r   &    61,845,869 & 6.2193\% \\
523.xalancbmk\_r &    18,381,250 & 4.6427\% \\
541.leela\_r     & 1,158,236,271 & 4.4438\% \\
519.lbm\_r       &   274,130,440 & 4.2130\% \\
531.deepsjeng\_r & 1,928,903,656 & 4.1749\% \\
525.x264\_r      &   767,981,981 & \textbf{3.2812\%} \\
544.nab\_r       &   151,310,616 & 2.2612\% \\
538.imagick\_r   &     2,424,899 & 2.0601\% \\
557.xz\_r        &    38,411,428 & 1.9069\% \\
505.mcf\_r       &   446,113,136 & 1.8263\% \\
508.namd\_r      &   327,718,955 & 1.0269\% \\
\bottomrule
\end{tabular}
\end{table}

\subsection{Store vs.\ Load Balance}

The load count (4.84B) is $2.01\times$ higher than the store count
(2.40B). This near-even 2:1 ratio reflects the H.264 decode pipeline's
balanced read-write pattern: for each macroblock, the decoder reads
residual coefficients and motion vectors from the bitstream buffers
(loads) and writes the reconstructed pixel block to the frame buffer
(stores). CABAC decoding reads probability tables (loads) and updates
them in place (stores). The deblocking filter similarly reads boundary
pixels and writes filtered values. The 2:1 ratio is structurally
characteristic of entropy-decoding pipelines that maintain
continuously updated probability or state tables alongside their
output buffers.

\subsection{Spill Fraction of Loads}

Within the ROI, the ratio of spills to total loads is:
\[
\frac{767{,}981{,}981}{4{,}837{,}558{,}326} \approx 15.9\%
\]
Approximately 1 in 6 load instructions is a spill reload. The JM
decoder's decode pipeline propagates a macroblock context structure
through multiple function layers: CABAC symbol decoding, inverse
quantisation, IDCT, and motion compensation each receive the same
macroblock descriptor, creating a set of live pointers (current
macroblock, reference frame list, slice header, SPS/PPS parameters)
that RISC-V's 32 registers cannot accommodate simultaneously across
all nested calls. The resulting 15.9\% spill-to-load fraction is
consistent with a moderately deep but non-recursive call chain.

\subsection{Pre-ROI Spill Density}

The pre-ROI spill rate of 3.31\% ($8{,}865$ spills over $268{,}194$
instructions) is consistent with the ROI rate (3.28\%), confirming
that the JM decoder's initialisation code shares the same structural
call patterns as the decode loop itself. The pre-ROI phase is
negligibly small (268,194 instructions), representing essentially
instantaneous startup overhead.

\subsection{ROI Coverage}

The ROI accounts for $23{,}405{,}758{,}981$ out of $23{,}406{,}228{,}013$
total simulated instructions:
\[
\frac{23{,}405{,}758{,}981}{23{,}406{,}228{,}013} \approx 99.998\%
\]
This is the highest ROI coverage of all evaluated benchmarks: the
decode loop begins almost immediately after process start, and the
JM decoder performs virtually no computation outside the ROI.

\subsection{H.264 Decode Pipeline Register Pressure}

The moderate 3.28\% spill rate reflects the mixed nature of H.264
decoding. The pipeline alternates between two structurally different
phases:
\begin{itemize}
    \item \textbf{Entropy decoding (CABAC/CAVLC)}: a symbol-by-symbol
          loop that maintains a compact state (range, offset, probability
          tables), fitting mostly within RISC-V's register file. Spill
          pressure is low in this phase.
    \item \textbf{Reconstruction (IDCT + motion compensation)}: a
          block-level operation that simultaneously requires: the residual
          coefficient array, the reference frame buffer pointer, the motion
          vector, the reconstructed macroblock buffer, and loop-filter
          boundary values. This phase generates higher register pressure,
          pulling the overall spill rate upward.
\end{itemize}
The resulting 3.28\% spill rate represents the time-weighted average
across these two phases. Compared to object-oriented workloads with
deep virtual dispatch chains (\texttt{omnetpp\_r} at 7.20\%,
\texttt{xalancbmk\_r} at 4.64\%), the H.264 decoder's more regular
data-flow pattern and shallower call depth yield a substantially lower
spill rate while still exhibiting non-trivial register pressure due
to the macroblock context propagation.

% =========================================================
\section{Conclusion}

The \texttt{525.x264\_r} benchmark (ldecod\_r) exhibits a register
spill rate of 3.28\% within the H.264 decode loop ROI on RISC-V
RV64GC, the eleventh highest across all 16 evaluated SPEC CPU2017
benchmarks. With 768 million spill events across 23.4 billion ROI
instructions and a near-perfect 99.998\% ROI coverage, the benchmark
is comprehensively characterised by this run. The 2.01:1 load-to-store
ratio reflects the balanced entropy-decode and reconstruction pipeline
of the JM reference decoder, and the 15.9\% spill-to-load fraction
arises from the macroblock context structure propagated through the
multi-layer decode call chain. These results place \texttt{x264\_r}
in the moderate register-pressure tier of the SPEC CPU2017 suite,
below the object-oriented and recursive workloads but above the
sequential data-processing and numerical benchmarks.

\end{document}
